\scanpage{50}
% archive sheet number: 49
This could be given in an extra paragraph (which I did not really touch
upon in the talk), containing also the deduction of the Weil conjectures
from the conjectures $C$ and $A$.

A last and rather trivial remark is the following. Let's introduce
\deleted{\unsure{a|$x$}} variants $A'_{\unsure{\ell|l|\chi}}(X)$ and $A''_{\ell}(X)$ as follows:

$A'_{\ell}(X)$ : if $2i\handwritten{\leq} n-1$, any element \inserted{$x$} of $H^{i}(X)$ whose image in $H^{i}(Y)$ is alge\handwritten{\unsure{(|\,}}braic, is algebraic.

$A''_{\ell}(X)$ : if $2\deleted{i}\unsure{j|i}\handwritten{\geq} n-1$, any \inserted{algebraic} element \inserted{\deleted{\unsure{$x$}}}\deleted{\unsure{of} $H^{\unsure{i|j}}(Y)$ whose image \unsure{in|on}}\handwritten{of}
$H^{\deleted{i}j+2}(X)$ is \deleted{algebraic, is algebraic.}\deleted{\handwritten{\unsure{is}}} the image of an algebraic element of $H^{i}(Y)$.

Let us consider also the specifications $A'_{\ell}(X)^{\unsure{o|0}}$ and $A''_{\ell}(X)^{o}$, where we
restrict to the \unsure{crittval|critival|critiwal} dimensions $2i=n-1$ if $n$ odd, $2i=n-2$ if $n$
even. All these conditions are in the nature of ``week'' Lefschets relations,
and they are trivially implied by $A_{\ell}(X)$ resp. $C_{\ell}(Y)$ (in the first
case, applying $\phi$ we see that $L_{X}x$ is algebraic; in the second, we take
$y=\handwritten{\Lambda}_{Y}\phi^{\handwritten{*}}(x)$). The remark then is that these pretend\inserted{\handwritten{ed}}ly ``week'' variants
in fact imply the full Lefs\deleted{\unsure{b}}chets relations for al\inserted{\handwritten{ge}}braic cycles, namely:

\uline{Proposition}\quad $\deleted{\unsure{A}}\unsure{C|E}_{\ell}(X)$ is equivalent to the conj\deleted{\unsure{an}}unction $C_{\ell}(Y)+A'_{\ell}(X\mathbin{\mathrm{x}}X)^{o}+{}$
$A''_{\ell}(X\mathbin{\mathrm{x}}X)^{o}$, hence (by induction) also to the conj\deleted{\unsure{an}}unction of the
conditions $ {A'_{\ell}}^{o}$ and $ {A''_{\ell}}^{o}$ for all of the varieties $X\mathbin{\mathrm{x}}X,Y\mathbin{\mathrm{x}}Y,Z\mathbin{\mathrm{x}}Z$ .... . Analoguous statement with $X\mathbin{\mathrm{x}}Y,Y\mathbin{\mathrm{x}}Z$ etc instead of \uline{$X\mathbin{\mathrm{x}}X,Y\mathbin{\mathrm{x}}Y$ etc.}

This comes from the remark that \deleted{if $X$ is even dimensional (say a
square like \unsure{$X\mathbin{\mathrm{x}}X$}) then} $A_{\ell}(X\deleted{\unsure{\mathbin{\mathrm{x}}X}})^{o}$ follows from the conjuction of
$A'_{\ell}(X)^{o}$ and $A''_{\ell}(X)^{o}$, as one sees by decomposing $L_{X}^{2}:H^{2m-2}(X)\handwritten{\to}H^{2m+2}(X)$
into $H^{2m-2}(X)\handwritten{\xrightarrow{\varphi^{*}}}H^{2m-2}(Y)\handwritten{\xrightarrow{\varphi_{*}}}H^{2m}(X)\handwritten{\overset{L_{X}}{-}}H^{2m+2}\deleted{\unsure{\mathrm{XX}}}(X)$ if $\dim X=2m$ is even,
and $H^{2m\deleted{\unsure{(}}+1-1}(X)\handwritten{\to}H^{2m+1+1}(X)$ into $H^{2m}(X)\handwritten{\overset{\varphi^{*}}{-}}H^{2m}(Y)\handwritten{\overset{\varphi_{*}}{\unsure{-|\to}}}H^{2m+2}(X)$ if
$\dim X=2m+1$ is odd.

Sincerely your's\\

